\documentclass[11pt]{article}
\usepackage[margin=0.75in]{geometry}
\usepackage{amsmath,amssymb,amsthm}
\usepackage{array,booktabs,xcolor,colortbl}
\usepackage{tikz}
\usepackage{hyperref}
\setlength{\parindent}{0pt}
\newcommand{\stepnum}[1]{%
  \tikz[baseline=(n.base)] \node[draw,circle,inner sep=1.6pt,
    font=\scriptsize\bfseries,fill=white] (n) {#1};}%
\newtheorem{theorem}{Theorem}
\theoremstyle{definition}
\newtheorem{definition}{Definition}
\title{Flow Proof Artifact: full\_adder}
\author{Generated by \texttt{flow doc proof}}
\date{\today}
\begin{document}
\maketitle
The full adder output matches binary addition with carry.\\[0.5em]
\textbf{Source.} patterson — Patterson \& Hennessy, *Computer Organization and Design*, §A.5\\[0.5em]
\bigskip
\noindent{\large \textbf{Derived fact 1.} \textit{Correct} \hfill \texttt{FullAdder/out.correct}}\par
\smallskip\noindent\textit{Coordinate: FullAdder · output · output matches specification}\par
\medskip
\noindent\fcolorbox{black!8}{%
\begin{minipage}{\dimexpr\textwidth-2\fboxsep-2\fboxrule}
\textbf{Goal.} We're showing that result.Sum equals expected.the conjunction of sum and result.Cout equals expected.carry.\\[0.35em]
$\displaystyle \forall A \in bit,\; result.Sum = expected.sum \land result.Cout = expected.carry$
\end{minipage}%
}\par
\medskip
\noindent
\begin{tabular}{@{} c >{\raggedright\arraybackslash}p{0.40\textwidth} @{\hspace{0.4em}\color{gray!50}\vrule width 0.5pt\hspace{0.4em}} c >{\raggedleft\arraybackslash}p{0.40\textwidth} @{}}
\multicolumn{2}{l}{\textbf{Proof}} & \multicolumn{2}{r}{\textbf{Mathematics}} \\
\midrule
\stepnum{1} & Let result = FullAdder(A, B, Cin). & & \multicolumn{1}{p{0.40\textwidth}}{} \\[0.55em]
\stepnum{2} & Let expected = binary\_add\_1bit(A, B, Cin). & & \multicolumn{1}{p{0.40\textwidth}}{} \\[0.55em]
\stepnum{3} & From ① and ②, this implies result.Sum equals expected.the conjunction of sum and result.Cout equals expected.carry. Hence proven. & \stepnum{3} & $result.Sum = expected.sum \land result.Cout = expected.carry$ \\[0.55em]
\bottomrule
\end{tabular}
\label{thm:FullAdder-output-output_matches_specification}
\par\medskip\hrule
\bigskip
\noindent{\large \textbf{Derived fact 2.} \textit{Correct} \hfill \texttt{Ripple4/out.correct}}\par
\smallskip\noindent\textit{Coordinate: Ripple4 · output · output matches specification}\par
\medskip
\noindent\fcolorbox{black!8}{%
\begin{minipage}{\dimexpr\textwidth-2\fboxsep-2\fboxrule}
\textbf{Goal.} We're showing that sum equals expected.\\[0.35em]
$\displaystyle \forall A \in bit,\; sum = expected$
\end{minipage}%
}\par
\medskip
\noindent
\begin{tabular}{@{} c >{\raggedright\arraybackslash}p{0.40\textwidth} @{\hspace{0.4em}\color{gray!50}\vrule width 0.5pt\hspace{0.4em}} c >{\raggedleft\arraybackslash}p{0.40\textwidth} @{}}
\multicolumn{2}{l}{\textbf{Proof}} & \multicolumn{2}{r}{\textbf{Mathematics}} \\
\midrule
\stepnum{1} & Let carry = Cin. & & \multicolumn{1}{p{0.40\textwidth}}{} \\[0.55em]
\stepnum{2} & Let step = FullAdder(A[i], B[i], carry). & & \multicolumn{1}{p{0.40\textwidth}}{} \\[0.55em]
\stepnum{3} & We invoke the derived fact governing output on FullAdder: the law that output matches specification, for output on FullAdder (instantiated for A[i], B[i], carry). & & \multicolumn{1}{p{0.40\textwidth}}{} \\[0.55em]
\stepnum{4} & Let sum = bits\_to\_int(s) + bits\_to\_int(carry) * 16. & & \multicolumn{1}{p{0.40\textwidth}}{} \\[0.55em]
\stepnum{5} & Let expected = bits\_to\_int(A) + bits\_to\_int(B) + int(Cin). & & \multicolumn{1}{p{0.40\textwidth}}{} \\[0.55em]
\stepnum{6} & From ①, ②, ③, ④, and ⑤, this implies sum equals expected. Hence proven. & \stepnum{6} & $sum = expected$ \\[0.55em]
\bottomrule
\end{tabular}
\label{thm:Ripple4-output-output_matches_specification}
\par\medskip\hrule
\end{document}
